
\documentclass{article}

\usepackage[left=3cm,top=3cm,right=2cm,bottom=2cm,a4paper]{geometry}
\usepackage{anyfontsize}
\usepackage[indent=0cm,skip=.1cm]{parskip}
\usepackage[brazil]{babel}

\begin{document}
\fontsize{12pt}{18pt}\selectfont

%=====================================
% CAPA
%=====================================

\begin{titlepage}

\begin{center}

FACULDADE MAURÍCIO DE NASSAU

\vspace{2cm}

Phellipe de Amorim martins - 01782567

Anthony Gabriel Lemos Baracho - 01813891

Ruan Kelvin Vieira dos Santos - 01797916

Everton Hasabias Furtunato Celestino - 1796727

Adriano Silva do Nascimento - 01769583

Fernanda Nogueira de França - 01735985

\vfill

{\fontsize{16pt}{20pt}\selectfont\textbf{SISTEMAS OPERACIONAIS}}

\vspace{0.5cm}


\vfill

Maceió - AL

Junho de 2026

\end{center}

\end{titlepage}

%=====================================
% FOLHA DE ROSTO
%=====================================

\begin{titlepage}

\begin{center}

Phellipe de Amorim martins - 01782567

Anthony Gabriel Lemos Baracho - 01813891

Ruan Kelvin Vieira dos Santos - 01797916

Everton Hasabias Furtunato Celestino - 1796727

Adriano Silva do Nascimento - 01769583

Fernanda Nogueira de França - 01735985

\vspace{6cm}

{\fontsize{16pt}{20pt}\selectfont\textbf{SISTEMAS OPERACIONAIS}}

\vspace{0.5cm}


\end{center}

\vspace{3cm}

\hspace*{7cm}
\parbox{8cm}{
Trabalho apresentado à disciplina de
Sistemas Operacionais do curso de
Ciências da Computação da Faculdade Maurício de Nassau,
como requisito para obtenção de nota.

\vspace{0.5cm}

Professor(a): Anderlan Henrique Batista Siqueira
}

\vfill

\begin{center}

Maceió - AL

Junho de 2026

\end{center}

\end{titlepage}

\tableofcontents
\newpag

\clearpage

\section{Introdução ao Conceito de Sistema Operacional}
Antes de compreender o que é o Linux, é fundamental estabelecer a definição de um Sistema Operacional (SO). Um SO é o software responsável por:
\begin{itemize}
    \item Gerenciar a parte física do computador (hardware), incluindo o processador, a memória RAM e os dispositivos de armazenamento.
    \item Coordenar e gerir a execução de outros programas e processos.
    \item Prover, em muitos casos, uma interface para a interação humana.
\end{itemize}
Sistemas como Windows da Microsoft, macOS da Apple e o próprio Linux são popularmente categorizados como sistemas operacionais devido ao cumprimento dessas funções.

\subsection{A Real Definição: O Linux como Kernel}
O principal mal-entendido entre o público geral é a concepção de que o Linux é um sistema operacional completo e pronto por si só. Tecnicamente, o \textbf{Linux é um Kernel} (núcleo).

O Kernel desempenha o papel mais crítico em um ecossistema computacional:
\begin{itemize}
    \item Atua como uma ponte de comunicação direta que traduz as instruções dos softwares para o hardware.
    \item Concentra e gerencia os drivers de todos os dispositivos conectados.
\end{itemize}
Diferente do Windows (que utiliza o Kernel \textit{NT}) e do macOS (que se baseia no Kernel \textit{Darwin}), o Linux não possui um kernel à parte; ele \textit{é} o próprio Kernel.

\subsection{A Analogia Automotiva}
Para fins de simplificação didática, o ecossistema dos sistemas operacionais pode ser comparado à indústria automobilística:
\begin{description}
    \item[Windows:] Representa um carro popular completo (com motor, rodas e carroceria) construído integralmente pela Microsoft, que detém o controle total e cobra por sua licença de uso.
    \item[macOS:] Equivale a um carro gratuito fornecido pela Apple, com a restrição de que o usuário só pode dirigir dentro de um ecossistema fechado (uma ``cidade'' controlada pela marca), onde o custo de vida é elevado.
    \item[Linux:] Não representa o veículo completo, mas sim o \textbf{motor}. Trata-se de um propulsor altamente versátil, gratuito e de código aberto. O usuário recebe o motor junto com o manual completo de instruções, possuindo total liberdade para utilizá-lo na construção de um carro, de uma moto ou de um gerador de energia elétrica.
\end{description}

\subsection{Distribuições Linux (Distros)}
Dado que um motor de maneira isolada não possui utilidade prática de locomoção, surgem as chamadas \textbf{Distribuições Linux}, ou simplesmente \textbf{Distros}. Elas são os veículos completos montados por comunidades ou corporações que unem o Kernel Linux a um conjunto de softwares e interfaces gráficas.
\begin{itemize}
    \item \textbf{Ubuntu:} Sistema desenvolvido pela empresa britânica Canonical, focado em entregar uma experiência familiar e completa para computadores de mesa e laptops.
    \item \textbf{Android:} Desenvolvido pela Google utilizando o Kernel Linux como base, adaptado para operar em smartphones de forma distinta dos desktops tradicionais.
    \item \textbf{ChromeOS:} Outra solução da Google baseada no mesmo núcleo, mas focada no mercado de laptops leves (Chromebooks).
\end{itemize}

\subsection{Desenvolvimento Colaborativo e Modelo Open Source}
O Kernel Linux foi originalmente concebido por Linus Torvalds em 1991, que ainda permanece como um de seus engenheiros principais. Devido à sua natureza de código aberto e licença permissiva para modificação, ele é mantido por uma vasta comunidade global:
\begin{itemize}
    \item Milhares de desenvolvedores voluntários e corporações líderes de tecnologia (incluindo a própria Microsoft) colaboram para seu aprimoramento.
    \item Cada entidade desenvolve e otimiza o código para atender às suas necessidades de mercado específicas, o que contribui de forma agregada para a robustez e versatilidade do kernel para todos os usuários.
\end{itemize}

\subsection{Presença do Linux no Mercado}
Apesar de não deter a liderança de mercado em computadores domésticos tradicionais (desktops), o Linux apresenta dominância absoluta em quase todas as outras verticais tecnológicas do planeta:
\begin{itemize}
    \item \textbf{Infraestrutura Web:} Servidores de internet e grandes centrais de processamento de dados (como data centers de hospedagem).
    \item \textbf{Dispositivos Móveis e IoT:} Bilhões de smartphones com sistema Android, televisores inteligentes e eletrodomésticos conectados.
    \item \textbf{Alta Tecnologia e Ciência:} Supercomputadores, robótica aeroespacial, caixas de autoatendimento bancário e sistemas de grande escala científica, como o acelerador de partículas do CERN.
\end{itemize}



\section{História do Linux}
A história do Linux começou em 1991, quando um estudante finlandês chamado Linus Torvalds iniciou um projeto pessoal com o objetivo de criar seu próprio núcleo de sistema operacional.
O que começou como um hobby acabou se tornando um dos maiores projetos colaborativos da história da computação.

Na época, Torvalds utilizava o sistema MINIX, criado para fins educacionais, mas queria explorar melhor os recursos do seu novo computador equipado com um processador Intel 80386. Assim, ele começou a desenvolver um emulador de terminal que evoluiu para um novo kernel.
O desenvolvimento inicial foi feito utilizando o compilador GNU C Compiler (GCC), que até hoje é amplamente utilizado na compilação do Linux.

Inicialmente, Torvalds pensou em chamar seu projeto de \textit{Freax}, uma combinação das palavras ``free'', ``freak'' e a letra ``x'', em referência ao Unix. No entanto, em setembro de 1991, Ari Lemmke, administrador do servidor da Universidade de Helsinki, decidiu disponibilizar o projeto sob o nome \textit{Linux}, sem consultar Torvalds.
Posteriormente, ele aceitou a ideia, e esse passou a ser o nome oficial.

Com o passar dos anos, diversos desenvolvedores passaram a contribuir com o projeto. Entre eles destacam-se Alan Cox, Marcelo Tosatti e Andrew Morton, que atuaram como mantenedores de diferentes versões do kernel. Entretanto, a maior parte do desenvolvimento sempre foi realizada pela comunidade de software livre, formada por milhares de programadores espalhados pelo mundo.
Além da colaboração da comunidade, diversas empresas também passaram a investir no desenvolvimento do Linux e de ferramentas complementares, tornando-o cada vez mais robusto e confiável.

Assim, o Linux nasceu como um projeto pessoal de um estudante universitário e, graças ao desenvolvimento colaborativo e ao modelo de software livre, tornou-se um dos sistemas operacionais mais importantes do mundo, sendo utilizado em servidores, supercomputadores, dispositivos embarcados e até mesmo no Android.



\section{Como o Linux Funciona Internamente}
Antes de tudo, é importante entender que o Linux é formado por várias camadas que trabalham juntas para permitir a comunicação entre o usuário e o hardware do computador.

Podemos imaginar essa estrutura como uma pirâmide. Na base está o hardware, acima dele fica o kernel, depois vêm as bibliotecas do sistema, o shell e, por último, as aplicações utilizadas pelo usuário.

Quando uma pessoa utiliza um programa, como um navegador ou um editor de texto, esse programa não conversa diretamente com o hardware. Toda comunicação passa pelo kernel, que é o núcleo do sistema operacional.

\subsection{Processo de Inicialização (Boot)}

O funcionamento do Linux começa no momento em que o computador é ligado.

A primeira coisa executada é a BIOS ou a UEFI, um pequeno software gravado na placa-mãe. Sua função é verificar se os componentes físicos do computador estão funcionando corretamente, como a memória RAM, o processador e os dispositivos conectados.

Depois dessa verificação, a BIOS ou UEFI procura um sistema operacional para iniciar. Nesse momento entra em ação o \textit{Bootloader}, que normalmente é o GRUB.

O GRUB é responsável por localizar o kernel Linux no disco rígido ou SSD e carregá-lo para a memória RAM.

A sequência de inicialização pode ser representada da seguinte forma:

\begin{center}
Computador $\rightarrow$ BIOS/UEFI $\rightarrow$ GRUB $\rightarrow$ Kernel Linux
\end{center}

Após ser carregado, o kernel assume completamente o controle do computador.

\subsection{O Kernel}

O kernel é a parte mais importante do Linux. Ele funciona como um intermediário entre os programas e o hardware.

Sempre que um aplicativo precisa realizar alguma tarefa, como abrir um arquivo, acessar a internet ou imprimir um documento, ele envia uma solicitação ao kernel.

O kernel recebe essa solicitação, verifica se ela pode ser executada, acessa o hardware necessário e devolve a resposta ao programa.

Por esse motivo, o kernel é frequentemente chamado de ``coração'' do sistema operacional.

\subsection{Gerenciamento de Processos}

Uma das principais funções do kernel é controlar os processos.

Um processo é qualquer programa que esteja em execução. Por exemplo, se uma pessoa estiver utilizando um navegador, um editor de texto e um reprodutor de música ao mesmo tempo, cada um deles será tratado pelo kernel como um processo diferente.

Como existe apenas um processador, o kernel precisa decidir qual processo utilizará a CPU em cada instante. Esse trabalho é realizado pelo escalonador de processos, também chamado de \textit{Scheduler}.

O \textit{Scheduler} divide o tempo do processador em pequenas frações, alternando rapidamente entre os programas. Essa troca acontece milhares de vezes por segundo, criando a sensação de que todos os aplicativos estão funcionando simultaneamente.

Além disso, cada processo possui um identificador chamado PID (\textit{Process ID}), que serve para diferenciá-lo dos demais.

Nas distribuições Linux modernas, o primeiro processo criado pelo kernel possui PID 1 e geralmente corresponde ao \textit{systemd}, responsável por iniciar todos os outros serviços do sistema.

\subsection{Gerenciamento de Memória}

Outra função essencial do kernel é controlar a memória RAM.

Cada programa precisa de uma determinada quantidade de memória para funcionar, e o kernel garante que um processo não interfira na área de memória utilizada por outro.

O Linux utiliza um mecanismo chamado \textit{Memória Virtual}. Nesse sistema, cada programa acredita que possui uma área exclusiva de memória, mesmo que esteja compartilhando a RAM com outros processos.

Quando a memória RAM está quase cheia, o Linux pode utilizar uma área do disco chamada \textit{Swap}, que funciona como uma memória auxiliar.

Embora o \textit{Swap} seja mais lento que a RAM, ele ajuda a evitar travamentos quando muitos programas estão abertos.

\subsection{Sistema de Arquivos}

No Linux, praticamente tudo é tratado como um arquivo.

Além dos documentos tradicionais, dispositivos como discos, teclados, mouses e impressoras também são representados como arquivos especiais.

Toda a estrutura do sistema começa no diretório raiz, representado pelo símbolo \texttt{/}.

A partir dele surgem diversas pastas importantes:

\begin{itemize}
\item \texttt{/boot}: contém os arquivos necessários para a inicialização do sistema;
\item \texttt{/home}: armazena os arquivos pessoais dos usuários;
\item \texttt{/etc}: guarda arquivos de configuração;
\item \texttt{/dev}: representa os dispositivos de hardware;
\item \texttt{/proc}: fornece informações sobre os processos e sobre o kernel;
\item \texttt{/var}: armazena logs e arquivos que mudam constantemente.
\end{itemize}

\subsection{Drivers de Dispositivos}

Os drivers são pequenos programas que permitem a comunicação entre o kernel e os componentes físicos do computador.

Quando um usuário conecta um pendrive, por exemplo, o kernel detecta automaticamente o novo dispositivo, identifica qual driver deve ser utilizado e disponibiliza seu acesso ao sistema.

\subsection{Chamadas de Sistema (\textit{System Calls})}

Os programas não podem acessar diretamente o hardware.

Sempre que precisam realizar alguma operação, eles utilizam funções chamadas \textit{System Calls}.

Por exemplo, ao abrir um arquivo, o programa solicita essa operação ao kernel através de uma chamada de sistema.

O kernel verifica se o arquivo existe, se o usuário possui permissão para acessá-lo e em qual local do disco ele está armazenado.

Somente depois dessas verificações o conteúdo é enviado ao programa.

Esse mecanismo aumenta a segurança e impede que aplicativos acessem recursos de maneira descontrolada.

\subsection{Segurança e Permissões}

O Linux possui um sistema robusto de permissões.

Cada arquivo possui um proprietário, um grupo e três tipos básicos de autorização:

\begin{itemize}
\item Leitura;
\item Escrita;
\item Execução.
\end{itemize}

Isso impede que usuários comuns alterem arquivos críticos do sistema sem autorização administrativa.

Esse modelo de segurança é um dos motivos pelos quais o Linux é amplamente utilizado em servidores e ambientes corporativos.

\subsection{Kernel Monolítico}

O Linux utiliza uma arquitetura chamada \textit{Kernel Monolítico}.

Isso significa que seus principais serviços, como gerenciamento de processos, memória, sistema de arquivos, drivers e rede, ficam integrados dentro do próprio kernel.

A principal vantagem desse modelo é o alto desempenho, pois a comunicação entre os componentes acontece de forma muito rápida.

Por outro lado, um erro em um driver ou em alguma parte do kernel pode afetar todo o sistema.

\subsection{Exemplo Prático}

Para entender melhor, imagine que um usuário deseja abrir uma fotografia armazenada no computador.

\begin{enumerate}
\item O usuário clica duas vezes sobre a imagem;
\item O visualizador de imagens envia uma solicitação para as bibliotecas do sistema;
\item As bibliotecas realizam uma chamada de sistema para o kernel;
\item O kernel verifica as permissões do usuário;
\item O kernel localiza o arquivo no SSD;
\item O hardware realiza a leitura dos dados;
\item O kernel devolve as informações ao programa;
\item O programa exibe a fotografia na tela.
\end{enumerate}

Todo esse processo acontece em poucos milissegundos.

\subsection{Conclusão}

Podemos concluir que o Linux funciona como um grande administrador dos recursos do computador.

Seu kernel atua como uma ponte entre os aplicativos e o hardware, controlando a memória, o processador, os dispositivos, os arquivos e a segurança do sistema.

Graças a essa arquitetura em camadas e ao eficiente gerenciamento de recursos, o Linux tornou-se um dos sistemas operacionais mais utilizados do mundo, estando presente em computadores pessoais, servidores, supercomputadores e dispositivos móveis, como o Android.


\section{Kali Linux}

\subsection{O que é o Kali Linux?}

O Kali Linux é uma distribuição Linux especializada em testes de segurança, análise de vulnerabilidades, auditoria de redes e computação forense. Ele é baseado no Debian e foi desenvolvido para fornecer um ambiente completo contendo ferramentas voltadas para profissionais de segurança da informação.

Diferentemente de distribuições como Ubuntu ou Fedora, que são projetadas para uso geral, o Kali Linux foi criado especificamente para atividades relacionadas à segurança digital. Seu objetivo principal é permitir que especialistas avaliem a segurança de sistemas, redes e aplicações, identificando vulnerabilidades antes que elas possam ser exploradas por pessoas mal-intencionadas.

\subsection{História do Kali Linux}

O Kali Linux foi lançado em 2013 pela empresa Offensive Security. Ele surgiu como sucessor de uma distribuição chamada \textit{BackTrack}, que já era bastante popular entre profissionais da área de segurança. Os desenvolvedores decidiram reconstruir completamente o sistema utilizando como base o Debian, tornando-o mais estável, seguro e fácil de manter.

Desde então, o Kali Linux recebe atualizações frequentes, incorporando novas ferramentas e acompanhando a evolução das ameaças e tecnologias de segurança. Ao longo dos anos, tornou-se uma das distribuições mais utilizadas por especialistas em segurança cibernética, instituições de ensino e equipes de resposta a incidentes.

\subsection{Como o Kali Linux é Utilizado}

O Kali Linux é frequentemente utilizado durante auditorias de segurança. Imagine uma empresa que deseja verificar se seus sistemas estão protegidos contra invasões. Profissionais autorizados podem utilizar o Kali Linux para simular ataques controlados e identificar falhas antes que criminosos as descubram.

Esse processo é conhecido como \textit{teste de invasão} (\textit{Penetration Test}). O objetivo não é causar danos, mas sim encontrar vulnerabilidades para que elas possam ser corrigidas.

Além disso, o sistema é amplamente utilizado em laboratórios educacionais para o ensino de redes, sistemas operacionais e segurança da informação.



\subsection{Principais Características}

Uma das principais características do Kali Linux é a grande quantidade de ferramentas especializadas que já vêm instaladas ou disponíveis em seus repositórios oficiais. Essas ferramentas são utilizadas para diversas finalidades, como:

\begin{itemize}
\item Análise de redes;
\item Auditoria de sistemas;
\item Testes de aplicações web;
\item Análise forense digital;
\item Monitoramento de tráfego;
\item Avaliação de vulnerabilidades.
\end{itemize}

Outra característica importante é sua flexibilidade. O Kali Linux pode ser executado em computadores físicos, máquinas virtuais, dispositivos móveis compatíveis e até mesmo a partir de um pendrive, sem necessidade de instalação no disco rígido do equipamento.

\subsection{Ferramentas Disponíveis}

O Kali Linux se destaca por oferecer centenas de ferramentas especializadas. Essas ferramentas são organizadas em categorias para facilitar sua utilização. Algumas dessas categorias incluem:

\begin{itemize}
\item Coleta de informações;
\item Análise de vulnerabilidades;
\item Segurança de aplicações web;
\item Engenharia reversa;
\item Computação forense;
\item Análise de redes;
\item Testes de segurança sem fio;
\item Monitoramento de tráfego.
\end{itemize}

Essa ampla coleção de ferramentas transforma o Kali Linux em uma plataforma completa para profissionais da área de segurança da informação.

\subsection{Conclusão}

Graças à sua ampla coleção de ferramentas, à estabilidade herdada do Debian e ao suporte da comunidade, o Kali Linux tornou-se uma das plataformas mais importantes da área de cibersegurança.

Mais do que um sistema operacional, ele representa um ambiente completo para pesquisa, aprendizado e fortalecimento da segurança digital. Seu uso responsável contribui para a identificação e correção de vulnerabilidades, promovendo ambientes computacionais mais seguros e preparados para enfrentar os desafios da segurança da informação contemporânea.

\clearpage


\begin{thebibliography}{99}

\bibitem{torvalds2001}
TORVALDS, Linus; DIAMOND, David.
\textit{Just for Fun: The Story of an Accidental Revolutionary}.
New York: HarperBusiness, 2001.

\bibitem{love2010}
LOVE, Robert.
\textit{Linux Kernel Development}.
3. ed. Upper Saddle River: Addison-Wesley, 2010.

\bibitem{tanenbaum2014}
TANENBAUM, Andrew S.; BOS, Herbert.
\textit{Modern Operating Systems}.
4. ed. Upper Saddle River: Pearson Prentice Hall, 2014.

\bibitem{kernelorg}
THE LINUX FOUNDATION.
\textit{About the Linux Kernel}.
Disponível em: \url{https://www.kernel.org}.
Acesso em: jun. 2026.

\bibitem{canonical}
CANONICAL LTD.
\textit{Ubuntu Documentation}.
Disponível em: \url{https://ubuntu.com/server/docs}.
Acesso em: jun. 2026.

\bibitem{android}
GOOGLE LLC.
\textit{Android Open Source Project}.
Disponível em: \url{https://source.android.com}.
Acesso em: jun. 2026.

\bibitem{kali}
OFFENSIVE SECURITY.
\textit{Kali Linux Documentation}.
Disponível em: \url{https://www.kali.org/docs}.
Acesso em: jun. 2026.

\bibitem{ali2011}
ALI, Shakeel; HERIYANTO, Tedi.
\textit{BackTrack 4: Assuring Security by Penetration Testing}.
Birmingham: Packt Publishing, 2011.

\bibitem{silberschatz2015}
SILBERSCHATZ, Abraham; GALVIN, Peter B.; GAGNE, Greg.
\textit{Fundamentos de Sistemas Operacionais}.
9. ed. Rio de Janeiro: LTC, 2015.

\bibitem{deitel2005}
DEITEL, Harvey M.; DEITEL, Paul J.; CHOFFNES, David R.
\textit{Sistemas Operacionais}.
3. ed. São Paulo: Pearson Prentice Hall, 2005.

\end{thebibliography}


\end{document}